\documentclass[11pt]{article}
\usepackage[margin=1in]{geometry}
\usepackage{amsmath,amssymb,bm}

\title{Crude CES Model with Explicit Matrix Entries}
\author{}
\date{\today}

\begin{document}
\maketitle

\section{Goal}
This document defines a deliberately crude but fully explicit RLGC model.
The model keeps inductive and capacitive coupling, includes dielectric substrate loss, and uses direct $\mathcal{O}(N^2)$ matrix assembly formulas.

\section{Geometry and Variables}
Discretize metal into $N$ straight segments.
Each segment $i$ has midpoint $\bm{r}_i=(x_i,y_i,z_i)$, length $l_i$, width $w_i$, thickness $t_i$, and direction unit vector $\hat{\bm{l}}_i$.

Define center-to-center distance
\[
d_{ij}=\|\bm{r}_i-\bm{r}_j\|,\qquad i\neq j.
\]
For self terms, use equivalent radius
\[
a_i=0.2235\,(w_i+t_i).
\]

Substrate model: one effective dielectric with
\[
\epsilon^*(\omega)=\epsilon'-j\frac{\sigma_d}{\omega},
\qquad
\epsilon_{\mathrm{eff}}^*(\omega)=\frac{1+\epsilon^*(\omega)}{2}.
\]

\section{Resistance Matrix $\bm{R}(\omega)$}
Diagonal only:
\[
R_{ii}(\omega)=\rho_i\frac{l_i}{w_i t_i}\,\Re\!\left[\kappa_i\coth(\kappa_i)\right],
\qquad
\kappa_i=(1+j)\frac{t_i}{\delta_i(\omega)},
\]
\[
\delta_i(\omega)=\sqrt{\frac{2\rho_i}{\omega\mu_i}}.
\]
Set
\[
R_{ij}(\omega)=0\quad (i\neq j).
\]

\section{Inductance Matrix $\bm{L}$}
Use explicit approximate entries.

Self terms:
\[
L_{ii}=\frac{\mu_0}{2\pi}l_i\left[\ln\!\left(\frac{2l_i}{a_i}\right)-1\right].
\]

Mutual terms (straight-segment crude approximation):
\[
L_{ij}=\frac{\mu_0}{4\pi}(\hat{\bm{l}}_i\cdot\hat{\bm{l}}_j)
\frac{l_i l_j}{\sqrt{d_{ij}^2+\eta^2}},
\qquad i\neq j,
\]
where $\eta=\tfrac12(a_i+a_j)$ avoids singular behavior at very small spacing.

Symmetrize after fill:
\[
\bm{L}\leftarrow\frac{\bm{L}+\bm{L}^\mathsf{T}}{2}.
\]

\section{Potential Matrix and Capacitance}
Build potential matrix $\bm{P}(\omega)$ directly.

Self terms:
\[
P_{ii}(\omega)=\frac{1}{2\pi\epsilon_0\epsilon_{\mathrm{eff}}^*(\omega)\,l_i}
\ln\!\left(\frac{2l_i}{a_i}\right).
\]

Mutual terms:
\[
P_{ij}(\omega)=\frac{1}{4\pi\epsilon_0\epsilon_{\mathrm{eff}}^*(\omega)}
\frac{1}{\sqrt{d_{ij}^2+\eta^2}},
\qquad i\neq j.
\]

Then
\[
\bm{C}^*(\omega)=\bm{P}(\omega)^{-1}.
\]
Use
\[
\bm{C}(\omega)=\Re[\bm{C}^*(\omega)],
\qquad
\bm{G}_d(\omega)=-\omega\,\Im[\bm{C}^*(\omega)].
\]

\section{Admittance Matrix and Solve}
Define branch impedance
\[
\bm{Z}_{RL}(\omega)=\bm{R}(\omega)+j\omega\bm{L},
\qquad
\bm{Y}_{RL}(\omega)=\bm{Z}_{RL}(\omega)^{-1}.
\]

For the compact crude model (segment-domain solve), use
\[
\bm{Y}(\omega)=\bm{Y}_{RL}(\omega)+j\omega\bm{C}(\omega)+\bm{G}_d(\omega).
\]
Given current excitation $\bm{i}(\omega)$,
\[
\bm{v}(\omega)=\bm{Y}(\omega)^{-1}\bm{i}(\omega).
\]

Equivalent impedance matrix:
\[
\bm{Z}(\omega)=\bm{Y}(\omega)^{-1}.
\]

\section{Approximation Type Used in Each Block}
\subsection{Geometry Discretization}
\textbf{Approximation type:} piecewise-uniform segment model (low-order PEEC style).

Each conductor is replaced by straight segments with uniform current and charge per segment.
This is a first-order spatial discretization, so current crowding near corners and rapidly varying charge density are not resolved.

\subsection{Equivalent Radius $a_i$}
\textbf{Approximation type:} rectangular-to-round shape reduction.

The mapping $a_i=0.2235(w_i+t_i)$ converts a rectangular cross section into a single effective radius.
This preserves the dominant logarithmic dependence in self terms but not detailed edge-field behavior.

\subsection{Resistance Matrix $\bm{R}(\omega)$}
\textbf{Approximation type:} 1D conduction with slab skin-effect correction.

The diagonal formula assumes current flows along segment length and uses a local thickness-only skin-depth correction.
Off-diagonal terms are set to zero, which neglects proximity-effect coupling between nearby conductors.

\subsection{Inductance Matrix $\bm{L}$}
\textbf{Approximation type:} partial inductance with filament-like mutual kernel.

The self term is a logarithmic long-wire approximation.
The mutual term uses centerline distance and orientation dot product, replacing the full Neumann double integral by a closed-form surrogate.
The softening parameter $\eta$ is a numerical regularization that removes near-singularity and acts as a finite-size correction.

\subsection{Potential Matrix $\bm{P}(\omega)$}
\textbf{Approximation type:} quasi-static Coulomb kernel in a homogeneous effective medium.

The mutual term is $1/r$ with distance softening; the self term is logarithmic.
This avoids full BEM panel integration and avoids layered Green functions, so vertical field stratification is represented only through an effective permittivity.

\subsection{Dielectric Substrate Model}
\textbf{Approximation type:} single effective-medium complex permittivity.

Using
\[
\epsilon_{\mathrm{eff}}^*(\omega)=\frac{1+\epsilon^*(\omega)}{2}
\]
is a quasi-TEM mixing rule for air plus substrate.
It captures first-order dielectric loading and loss tangent behavior but not full multilayer dispersion or anisotropy.

\subsection{Capacitance and Dielectric Conductance Split}
\textbf{Approximation type:} complex-capacitance decomposition.

The model computes $\bm{C}^*=\bm{P}^{-1}$ and then separates storage and loss as
$\bm{C}=\Re(\bm{C}^*)$ and $\bm{G}_d=-\omega\Im(\bm{C}^*)$.
This is a linear small-signal representation and assumes weak field nonlinearity.

\subsection{System Solve}
\textbf{Approximation type:} linear frequency-domain lumped network.

The final admittance
\[
\bm{Y}=\bm{Y}_{RL}+j\omega\bm{C}+\bm{G}_d
\]
assumes linear, time-harmonic behavior and no radiation coupling.
Hence it is most reliable in electrically compact, quasi-static or moderately low-frequency regimes.


\end{document}
